
\subsection{从 $A_\infty$ 空间到拓扑幺半群的
$A_\infty$ 映射与 $\mathcal N'_{k+1}$}
\label{subsec:7-5-6}

本小节证明命题~\ref{prop:7-5-43}，
并说明 $\mathcal N'_{k+1}$ 与 $\mathcal N_{k+1}$ 的关系。
命题~\ref{prop:7-5-43} 是斯塔谢夫已知的结果
\textbf{[Sta70]}。事实上，可以取
\begin{equation}
\mathcal N'_{k+1}=\mathcal M_{k+2}^{\mathrm{main}}.
\tag{7.5.59}\label{eq:7-5-59}
\end{equation}
为了证明命题~\ref{prop:7-5-43}，
只须证明式~\eqref{eq:7-5-59} 的边界
由式~\eqref{eq:7-5-42-1}、\eqref{eq:7-5-42-2}
所描述；这是因为我们已在第 2.1.2 小节证明
$\mathcal M_{k+2}^{\mathrm{main}}$ 同胚于 $D^{k-1}$。

令 $(\Sigma,\vec z)\in
\partial\mathcal M_{k+2}^{\mathrm{main}}$ 为一个边界点，并写
\[
\vec z=(z_0,z_1,\ldots,z_k,z_{k+1})
\in(\partial\Sigma)^{k+2}.
\]
为了考察 $\mathcal M_{k+2}^{\mathrm{main}}$ 的
余维一边界，只须考虑 $\Sigma$ 恰有两个分支
$\Sigma_0$、$\Sigma_1$ 的情形。
不妨设 $z_0\in\Sigma_0$，并置
\[
\{z_{\mathrm{sing}}\}=\Sigma_0\cap\Sigma_1.
\]
分成以下两种情形：
\begin{align}
z_{k+1}&\in\partial\Sigma_0
&&\text{（图~\ref{fig:7-5-2}）},
\tag{7.5.60}\label{eq:7-5-60}\\
z_{k+1}&\in\partial\Sigma_1
&&\text{（图~\ref{fig:7-5-3}）}.
\tag{7.5.61}\label{eq:7-5-61}
\end{align}

\begin{figure}[!htbp]
 \centering
 \setlength{\unitlength}{0.78mm}
 \begin{picture}(122,60)
  \linethickness{0.48pt}
  \put(44,28){\circle{32}}
  \put(76,28){\circle{32}}
  \linethickness{0.32pt}
  \Line(27,37)(29,42) \Line(31,45)(36,48)
  \Line(39,49)(45,50) \Line(48,50)(53,48)
  \Line(56,47)(59,44)
  \Line(27,34)(26,29) \Line(26,26)(27,21)
  \Line(28,18)(31,14) \Line(33,12)(38,9)
  \Line(41,8)(47,8) \Line(50,9)(55,11)
  \Line(58,13)(60,16)
  \Line(64,44)(68,47) \Line(71,49)(77,50)
  \Line(80,50)(85,48) \Line(88,46)(91,42)
  \Line(93,39)(94,35)
  \Line(94,21)(92,17) \Line(90,14)(86,11)
  \Line(83,9)(77,8) \Line(74,8)(69,10)
  \Line(66,12)(63,16)
  \linethickness{0.48pt}
  \put(60,28){\circle*{1.5}}
  \put(30,38){\makebox(0,0){\scriptsize$\times$}}
  \put(44,44){\makebox(0,0){\scriptsize$\times$}}
  \put(54,40){\makebox(0,0){\scriptsize$\times$}}
  \put(29,29){\makebox(0,0){\scriptsize$\times$}}
  \put(34,16){\makebox(0,0){\scriptsize$\times$}}
  \put(50,13){\makebox(0,0){\scriptsize$\times$}}
  \put(56,18){\makebox(0,0){\scriptsize$\times$}}
  \put(65,39){\makebox(0,0){\scriptsize$\times$}}
  \put(87,39){\makebox(0,0){\scriptsize$\times$}}
  \put(92,28){\makebox(0,0){\scriptsize$\times$}}
  \put(88,17){\makebox(0,0){\scriptsize$\times$}}
  \put(44,28){\makebox(0,0){$\Sigma_1$}}
  \put(76,28){\makebox(0,0){$\Sigma_0$}}
  \put(66,56){\makebox(0,0){$z_{\mathrm{sing}}$}}
  \Line(64,52)(61,30)
  \put(55,43){\makebox(0,0)[l]{$z_{i+1}$}}
  \put(51,10){\makebox(0,0)[l]{$z_{i+\ell-1}$}}
  \put(57,16){\makebox(0,0)[l]{$z_{i+\ell}$}}
  \put(65,42){\makebox(0,0)[l]{$z_i$}}
  \put(89,41){\makebox(0,0)[l]{$z_1$}}
  \put(94,28){\makebox(0,0)[l]{$z_0$}}
  \put(90,16){\makebox(0,0)[l]{$z_{k+1}$}}
 \end{picture}
 \par\smallskip
 \renewcommand{\thefigure}{7.5.2}
 \refstepcounter{figure}
 \textbf{图~\thefigure}\label{fig:7-5-2}
\end{figure}

\begin{figure}[!htbp]
 \centering
 \setlength{\unitlength}{0.78mm}
 \begin{picture}(122,59)
  \linethickness{0.48pt}
  \put(44,27){\circle{32}}
  \put(76,27){\circle{32}}
  \linethickness{0.32pt}
  \Line(27,36)(29,41) \Line(31,44)(36,47)
  \Line(39,48)(45,49) \Line(48,49)(53,47)
  \Line(56,46)(59,43)
  \Line(27,33)(26,28) \Line(26,25)(27,20)
  \Line(28,17)(31,13)
  \Line(64,43)(68,46) \Line(71,48)(77,49)
  \Line(80,49)(85,47) \Line(88,45)(91,41)
  \Line(93,38)(94,34)
  \linethickness{0.48pt}
  \put(60,27){\circle*{1.5}}
  \put(54,39){\makebox(0,0){\scriptsize$\times$}}
  \put(65,38){\makebox(0,0){\scriptsize$\times$}}
  \put(34,15){\makebox(0,0){\scriptsize$\times$}}
  \put(52,17){\makebox(0,0){\scriptsize$\times$}}
  \put(92,27){\makebox(0,0){\scriptsize$\times$}}
  \put(44,27){\makebox(0,0){$\Sigma_1$}}
  \put(76,27){\makebox(0,0){$\Sigma_0$}}
  \put(66,55){\makebox(0,0){$z_{\mathrm{sing}}$}}
  \Line(64,51)(61,29)
  \put(54,42){\makebox(0,0)[l]{$z_{\ell+1}$}}
  \put(65,41){\makebox(0,0)[l]{$z_\ell$}}
  \put(31,12){\makebox(0,0)[r]{$z_k$}}
  \put(53,14){\makebox(0,0)[l]{$z_{k+1}$}}
  \put(94,27){\makebox(0,0)[l]{$z_0$}}
 \end{picture}
 \par\smallskip
 \renewcommand{\thefigure}{7.5.3}
 \refstepcounter{figure}
 \textbf{图~\thefigure}\label{fig:7-5-3}
\end{figure}

在情形~\eqref{eq:7-5-60} 中，设
\[
z_{\mathrm{sing}}\in[z_i,z_{i+\ell+1}],
\qquad
z_i,z_{i+\ell+1}\in\Sigma_0;
\]
见图~\ref{fig:7-5-2}。
注意 $i\ne k+1$，因为 $\Sigma_1$ 上至少应有一个标记点。
于是
\[
(\Sigma_0,
 (z_0,z_1,\ldots,z_i,z_{\mathrm{sing}},
  z_{i+\ell+1},\ldots,z_{k+1}))
\in\mathcal M_{k+2-\ell+1}^{\mathrm{main}}
=\mathcal N'_{k-\ell+2}
\]
且
\[
(\Sigma_1,(z_{\mathrm{sing}},z_{i+1},\ldots,z_{i+\ell}))
\in\mathcal M_{\ell+1}^{\mathrm{main}}.
\]
也就是说，情形~\eqref{eq:7-5-60}
对应于式~\eqref{eq:7-5-42-1}。

下面考虑情形~\eqref{eq:7-5-61}。
存在 $\ell$，使得
$z_1,\ldots,z_\ell\in\Sigma_0$，
$z_{\ell+1},\ldots,z_{k+1}\in\Sigma_1$。
于是
\[
(\Sigma_0,(z_0,z_1,\ldots,z_\ell,z_{\mathrm{sing}}))
\in\mathcal M_{\ell+2}^{\mathrm{main}}
=\mathcal N'_{\ell+1}
\]
且
\[
(\Sigma_1,(z_{\mathrm{sing}},z_{\ell+1},\ldots,z_k,z_{k+1}))
\in\mathcal M_{k-\ell+2}^{\mathrm{main}}
=\mathcal N'_{k-\ell+1}.
\]
因此，情形~\eqref{eq:7-5-61}
对应于式~\eqref{eq:7-5-42-2}。
由此证明了式~\eqref{eq:7-5-59}
满足所要求的性质。\hfill\(\square\)

\begin{remark}{7.5.62}\label{rem:7-5-62}
斯塔谢夫在 \textbf{[Sta70]} 中使用
$\mathcal M_{k+2}^{\mathrm{main}}=\mathcal N'_{k+1}$
来定义从 $A_\infty$ 空间到幺半群的
$A_\infty$ 映射；这同我们在
定义~\ref{def:7-1-66} 中使用 $\mathcal N_{k+1}$
来定义两个 $A_\infty$ 空间之间的
$A_\infty$ 映射完全类似。
\end{remark}

现在证明下面的定理~\ref{thm:7-5-63}。
回忆，空间 $\mathcal N_{k+1}$ 上的等价关系
$\sim$（定义~\ref{def:7-1-54}）由式 (7.5.44) 定义。

\begin{theorem}{7.5.63}\label{thm:7-5-63}
$\mathcal N_{k+1}/\!\sim$ 同胚于
$\mathcal M_{k+2}^{\mathrm{main}}$。
\end{theorem}

\begin{proof}
我们将给出从 $\mathcal N_{k+1}/\!\sim$
到 $\mathcal M_{k+2}^{\mathrm{main}}$ 的显式同胚。
令
\[
S=(\Sigma,\vec z,\rho)\in\mathcal N_{k+1}.
\]
先考虑 $1\notin\operatorname{Im}(\rho)$ 的情形。
像第~\ref{subsec:4-6-1} 小节那样把 $\Sigma$ 分解为分支，
并设其各分支为
\[
(\Sigma_\alpha,\vec z_\alpha)
\in\mathcal M_{k(\alpha)+1}^{\mathrm{main}}
\qquad(\alpha\in\mathcal A).
\]
令 $\rho_\alpha$ 为赋给 $\Sigma_\alpha$ 的时间分配。
依定义~\ref{def:4-6-6} 在分支集上定义偏序。
设 $\Sigma_{\alpha_0}$ 为最大分支，
即含有第 $0$ 个标记点 $z_0$ 的分支。

使用第~\ref{sec:3-4} 节定义的同胚
\[
\overline\Theta\colon
\overline{\operatorname{Gr}}_{k(\alpha)+1}
\longrightarrow
\mathcal M_{k(\alpha)+1}^{\mathrm{main}};
\]
亦见 \textbf{[FuOh97]}。置
\[
(\Sigma_\alpha,\vec z_\alpha)
=\overline\Theta(t_\alpha,\ell_\alpha),
\qquad
(\Sigma,\vec z)=\overline\Theta(t,\ell).
\]
这里以及下文均使用第~\ref{sec:3-4} 节的记号。
注意，$t$ 是把各 $t_\alpha$ 的边连接起来得到的。
更准确地说，可以把
\[
C^1(t)=\bigcup_{\alpha\in\mathcal A}C^1(t_\alpha)
\]
以及
\[
C^1(t_\alpha)\cap C^1(t_{\alpha'})=
\begin{cases}
\text{一条边},&\Sigma_\alpha\cap\Sigma_{\alpha'}\ne\varnothing,\\
\varnothing,&\text{否则}
\end{cases}
\]
作为等同。
这里 $C^1(t)$ 是 $t$ 的边集。
若 $\Sigma_\alpha\cap\Sigma_{\alpha'}\ne\varnothing$
且 $\Sigma_\alpha\leq\Sigma_{\alpha'}$，
就把
$C^1(t_\alpha)\cap C^1(t_{\alpha'})$
中的边记为 $e_\alpha$。
对每个 $\alpha\ne\alpha_0$，
这样的 $e_\alpha$ 唯一存在；见图~\ref{fig:7-5-4}。
令 $e_{\alpha_0}$ 为 $t_{\alpha_0}$ 中
含有 $t$ 的第 $0$ 条外边的那条外边。
\translatornote{源书此处把 $t_{\alpha_0}$ 误写为未定义的 $t_0$；
译文按本段始终使用的分支指标及图~7.5.4 补回 $\alpha_0$。
见疑误 A-0630。}
于是
\begin{equation}
C^1_{\mathrm{int}}(t)
=\bigcup_{\alpha\in\mathcal A}C^1_{\mathrm{int}}(t_\alpha)
 \cup\{e_\alpha\mid
       \alpha\in\mathcal A,\ \alpha\ne\alpha_0\},
\tag{7.5.64}\label{eq:7-5-64}
\end{equation}
且式~\eqref{eq:7-5-64} 右端是不交并。

\begin{figure}[!htbp]
 \centering
 \setlength{\unitlength}{0.82mm}
 \begin{picture}(132,88)
  \linethickness{0.48pt}
  \Line(82,44)(61,65)
  \Line(61,65)(38,84)
  \Line(61,65)(28,65)
  \Line(61,65)(39,47)
  \Line(82,44)(62,24)
  \Line(62,24)(31,15)
  \Line(62,24)(51,2)
  \Line(82,44)(114,44)
  \put(114,44){\vector(1,0){9}}
  \put(59,67){\makebox(0,0)[r]{$\alpha_1$}}
  \put(38,86){\makebox(0,0)[r]{$\alpha_3$}}
  \put(27,67){\makebox(0,0)[r]{$\alpha_4$}}
  \put(39,45){\makebox(0,0)[r]{$\alpha_5$}}
  \put(60,22){\makebox(0,0)[r]{$\alpha_2$}}
  \put(30,13){\makebox(0,0)[r]{$\alpha_6$}}
  \put(51,0){\makebox(0,0)[t]{$\alpha_7$}}
  \put(47,78){\makebox(0,0)[l]{$\rho_{\alpha_3}$}}
  \put(38,68){\makebox(0,0)[l]{$\rho_{\alpha_4}$}}
  \put(47,50){\makebox(0,0)[l]{$\rho_{\alpha_5}$}}
  \put(67,58){\makebox(0,0)[l]{$\rho_{\alpha_1}$}}
  \put(68,29){\makebox(0,0)[l]{$\rho_{\alpha_2}$}}
  \put(42,19){\makebox(0,0)[l]{$\rho_{\alpha_6}$}}
  \put(55,10){\makebox(0,0)[l]{$\rho_{\alpha_7}$}}
  \put(92,47){\makebox(0,0)[b]{$\rho_{\alpha_0}$}}
  \put(126,49){\makebox(0,0)[l]{$\nu_1=$}}
  \put(126,41){\makebox(0,0)[l]{第 $0$ 条外边}}
 \end{picture}
 \par\smallskip
 \renewcommand{\thefigure}{7.5.4}
 \refstepcounter{figure}
 \textbf{图~\thefigure}\label{fig:7-5-4}
\end{figure}

注意 $\ell(e_\alpha)=\infty$；见第~\ref{subsec:7-1-3} 小节。
下面利用时间分配 $\rho_\alpha$ 的信息修改
$\ell$，或者说修改
\[
\ell_\alpha\colon
C^1_{\mathrm{int}}(t_\alpha)\longrightarrow[0,\infty],
\]
从而如下得到
\[
\ell'\colon C^1_{\mathrm{int}}(t)\longrightarrow[0,\infty].
\]

\begin{definition}{7.5.65}\label{def:7-5-65}
若 $e\in C^1_{\mathrm{int}}(t_\alpha)$，置
\begin{equation}
\ell'(e)=
\frac{\ell_\alpha(e)(1-\rho_\alpha)}
{\rho_\alpha(\rho_\alpha+\ell_\alpha(e))}.
\tag{7.5.66.1}\label{eq:7-5-66-1}
\end{equation}
对 $\alpha\in\mathcal A$、$\alpha\ne\alpha_0$，置
\begin{equation}
\ell'(e_\alpha)=\frac{1-\rho_\alpha}{\rho_{\alpha'}},
\tag{7.5.66.2}\label{eq:7-5-66-2}
\end{equation}
其中
$\Sigma_\alpha\cap\Sigma_{\alpha'}\ne\varnothing$、
$\Sigma_\alpha\leq\Sigma_{\alpha'}$，且
$e_\alpha=C^1(t_\alpha)\cap C^1(t_{\alpha'})$。
\end{definition}

这样便得到
\begin{equation}
F_0(\Sigma,\vec z,\rho)
=\overline\Theta(t,\ell')
\in\mathcal M_{k+1}^{\mathrm{main}}.
\tag{7.5.67}\label{eq:7-5-67}
\end{equation}
源书接着说，在定义
$F_0(\Sigma,\vec z,\rho)$ 时没有使用
$\rho_{\alpha_0}$；利用 $\rho_{\alpha_0}$
可再定义一个标记点，从而得到
$\mathcal M_{k+2}^{\mathrm{main}}$ 的元素
\translatornote{按定义~\ref{def:7-5-65} 的字面，
当 $e\in C^1_{\mathrm{int}}(t_{\alpha_0})$ 时，
式~\eqref{eq:7-5-66-1} 实际含有 $\rho_{\alpha_0}$；
而源书说要由它得到 $\mathcal M_{k+2}^{\mathrm{main}}$ 的元素后，
又只重复写出紧接着的 $\mathcal M_{k+1}^{\mathrm{main}}$ 元素。
后面的引理~\ref{lem:7-5-70} 明确把该式用作底空间映射，
再以 $\rho_{\alpha_0}\in[0,1)$ 给出纤维坐标并提升为 $F$；
译文保留公式和后续构造，不擅改其余依赖关系。
见疑误 A-0631。}。
先写
\[
F_0(\Sigma,\vec z,\rho)
=(\Sigma',\vec z\,')
\in\mathcal M_{k+1}^{\mathrm{main}}.
\]
用它证明下面的引理。置
\translatornote{源书式 (7.5.69) 写成 $z_0,z_{k+2}$；
但 $\mathcal M_{k+2}^{\mathrm{main}}$ 的 $k+2$ 个标记点
编号为 $z_0,\ldots,z_{k+1}$，且紧接的
$\operatorname{Forget}_{k+2}$ 正是遗忘最后一个、即第 $k+2$ 个
标记点 $z_{k+1}$。译文据此改正下标。见疑误 A-0632。}
\begin{align}
X&=\{(\Sigma,\vec z,\rho)\in\mathcal N_{k+1}
      \mid1\notin\operatorname{Im}(\rho)\},
\tag{7.5.68}\label{eq:7-5-68}\\
Y&=\{(\Sigma',\vec z\,')\in\mathcal M_{k+2}^{\mathrm{main}}
      \mid z_0,z_{k+1}\text{ 位于 }\Sigma'
      \text{ 的同一分支上}\}.
\tag{7.5.69}\label{eq:7-5-69}
\end{align}
令下列映射遗忘最后一个（第 $k+2$ 个）标记点：
\[
\operatorname{Forget}_{k+2}\colon
\mathcal M_{k+2}^{\mathrm{main}}
\longrightarrow\mathcal M_{k+1}^{\mathrm{main}}
\]
并以 $\overline Y$ 表示 $Y$ 的像。
注意，$Y\to\overline Y$ 是纤维同胚于 $[0,1)$ 的纤维丛。

\begin{lemma}{7.5.70}\label{lem:7-5-70}
映射 $F_0$ 定义了映射
$F_0\colon X\to\overline Y$，
它可提升为同胚 $F\colon X\to Y$。
\end{lemma}

\begin{proof}
把 $\overline Y\times[0,1)$ 与 $Y$ 等同。置
\[
F(\Sigma,\vec z,\rho)
=(\rho_{\alpha_0},F_0(\Sigma,\vec z,\rho))
\in[0,1)\times\overline Y.
\]
注意：
\translatornote{定义~\ref{def:7-1-53} 允许时间分配取 $0$，
而 $X$ 只排除了取值 $1$；式~\eqref{eq:7-5-66-1}、
\eqref{eq:7-5-66-2} 以及下面逆映射中的阈值却含有
$\rho_\alpha$、$\rho_{\alpha'}$ 或 $\rho_{\alpha_0}$ 作分母。
源书没有说明在零端点采用何种扩展值约定，
故这些公式按原样保留，零端点的连续延拓仍待核。
见疑误 A-0633。}
\begin{enumerate}[label=\textup{(7.5.71.\arabic*)},leftmargin=6em]
\item\label{eq:7-5-71-1}
对每个固定的 $\rho_\alpha$，
式~\eqref{eq:7-5-66-1} 所定义的函数
$\ell_\alpha(e)\mapsto\ell'(e)$ 单调递增，并且
\[
\lim_{\ell_\alpha(e)\to0}\ell'(e)=0,
\qquad
\lim_{\ell_\alpha(e)\to\infty}\ell'(e)
=\frac{1-\rho_\alpha}{\rho_\alpha}.
\]
\item\label{eq:7-5-71-2}
对每个固定的 $\rho_\alpha$，
式~\eqref{eq:7-5-66-2} 所定义的函数
$\rho_{\alpha'}\mapsto\ell'(e_\alpha)$ 单调递减，并且
\[
\lim_{\rho_{\alpha'}\to\rho_\alpha}\ell'(e_\alpha)
=\frac{1-\rho_\alpha}{\rho_\alpha},
\qquad
\lim_{\rho_{\alpha'}\to0}\ell'(e_\alpha)=\infty.
\]
\end{enumerate}
利用~\ref{eq:7-5-71-1}、\ref{eq:7-5-71-2}
可以证明 $F$ 的连续性。

下面定义其逆映射。
令
\[
(c,\overline\Theta(t,\ell'))
\in[0,1)\times\overline Y\cong Y,
\qquad
\rho_{\alpha_0}=c.
\]
设 $v_0\in C^0_{\mathrm{int}}(t)$ 是 $t$ 的最末内顶点。
考虑所有满足下述条件的顶点 $v$ 所成集合 $V_{\alpha_0}$：
$v$ 可由一条通向 $v_0$ 的路径连接，而该路径只由满足
\[
\ell'(e)<\frac{1-\rho_{\alpha_0}}{\rho_{\alpha_0}}
\]
的边 $e$ 组成。
令 $t_{\alpha_0}$ 为 $t$ 的连通子图，使得
\[
C^0_{\mathrm{int}}(t_{\alpha_0})=V_{\alpha_0}.
\]
记
$\ell'_{\alpha_0}
=\ell'|_{C^1_{\mathrm{int}}(t_{\alpha_0})}$，
并定义
\[
\ell_{\alpha_0}\colon
C^1_{\mathrm{int}}(t_{\alpha_0})\longrightarrow(0,\infty]
\]
使它满足方程
\[
\ell'_{\alpha_0}(e)
=\frac{\ell_{\alpha_0}(e)(1-\rho_{\alpha_0})}
{\rho_{\alpha_0}
 (\rho_{\alpha_0}+\ell_{\alpha_0}(e))}.
\]
由~\ref{eq:7-5-71-1}，这样的 $\ell_{\alpha_0}$ 唯一存在。

考虑补集 $t\setminus t_{\alpha_0}$。
设 $t_{\alpha_i}$（$i=1,2,\ldots$）
为 $t\setminus t_{\alpha_0}$ 的连通分支，
并令 $v_{\alpha_i}$ 为 $t_{\alpha_i}$ 中离 $v_0$ 最近的顶点。
唯一存在一条边 $e_{\alpha_i}$，
把 $v_{\alpha_i}$ 与 $t_{\alpha_0}$ 的一个顶点连接起来。
\translatornote{源书称该边连接到“$v_{\alpha_0}$ 的一个顶点”，
但 $v_{\alpha_0}$ 未定义且单个顶点不能再含顶点；
依紧前定义的连通子图及下一步分层构造，译文改为
$t_{\alpha_0}$ 的一个顶点。见疑误 A-0634。}
由下式定义 $\rho_{\alpha_i}$：
\[
\ell'(e_{\alpha_i})
=\frac{1-\rho_{\alpha_0}}{\rho_{\alpha_i}}.
\]
利用~\ref{eq:7-5-71-2} 以及由
$t_{\alpha_0}$、$t_{\alpha_i}$ 的定义直接得到的不等式
\[
\ell'(e_{\alpha_i})
\geq\frac{1-\rho_{\alpha_0}}{\rho_{\alpha_0}},
\]
容易看出
\[
\rho_{\alpha_i}\leq\rho_{\alpha_0},
\qquad i=1,2,\ldots.
\]
接着定义
\[
\ell_{\alpha_i}\colon
C^1_{\mathrm{int}}(t_{\alpha_i})\longrightarrow(0,\infty]
\]
使它满足方程
\[
\ell'_{\alpha_i}(e)
=\frac{\ell_{\alpha_i}(e)(1-\rho_{\alpha_i})}
{\rho_{\alpha_i}
 (\rho_{\alpha_i}+\ell_{\alpha_i}(e))},
\qquad
\ell'_{\alpha_i}
=\ell'|_{C^1_{\mathrm{int}}(t_{\alpha_i})}.
\]

然后取补集
\translatornote{源书先把全部 $t_{\alpha_i}$ 定义为
$t\setminus t_{\alpha_0}$ 的连通分支，随即又要求对
$t\setminus\bigcup_i t_{\alpha_i}$ 的连通分支重复构造；
按字面该补集只留下已处理的 $t_{\alpha_0}$，
递归层次因而并不清楚。译文保留源步骤，不擅改成猜测的迭代集合。
见疑误 A-0635。}
$t\setminus\bigcup_i t_{\alpha_i}$ 及其各连通分支，
以同样的方法定义
$t_{\alpha_j}$、$\rho_{\alpha_j}$ 和 $\ell_{\alpha_j}$。
由此从
$(c,\overline\Theta(t,\ell'))
\in[0,1)\times\overline Y\cong Y$
得到
\[
t=\bigcup_\alpha t_\alpha,
\qquad \rho_\alpha,
\qquad \ell_\alpha.
\]
最后由这些子树重建曲线，置
\[
(\Sigma_\alpha,\vec z_\alpha)
=\overline\Theta(t_\alpha,\ell_\alpha),
\qquad
(\Sigma,\vec z)=\overline\Theta(t,\ell).
\]
容易看出，这样得到的正是 $F$ 的逆映射。
\end{proof}

最后证明 $F$ 可延拓为同胚
\[
\mathcal N_{k+1}/\!\sim
\longrightarrow\mathcal M_{k+2}^{\mathrm{main}}.
\]
首先注意，当 $\rho_\alpha\to1$ 时，
式~\eqref{eq:7-5-66-1}、\eqref{eq:7-5-66-2}
的右端趋于 $0$。
这意味着时间分配为 $1$ 的部分在 $F$ 的像中收缩为一点。
而式 (7.5.44) 中的等价关系 $\sim$
正是通过收缩时间分配为 $1$ 的分支定义的。

还要注意，空间 $\mathcal N_{k+1}/\!\sim$
的边界（余维一层）也像
命题~\ref{prop:7-5-43} 那样分解。
利用以上两点，容易把 $F$ 延拓为所需的同胚。
\end{proof}
